% ===== §5 =====
\section{The Recognition of Shared Value}
\label{sec:sharedvalue}

\noindent
A second work of relational understanding is the recognition of shared value, the coming into view, between two people, of what they hold in common and can hold together. Here too there are established accounts, and they reach toward shared value from two directions.

One direction treats shared value as a similarity to be discovered. A body of research on close relationships finds that partners who share values report greater satisfaction and stability, and models the matter as a matching, on which two people bring antecedent value profiles and the relationship prospers to the degree that these profiles coincide~\citep{byrne1971attraction}. Shared value, on this rendering, is a pre-existing overlap that the partners find in one another. A second direction treats shared value as a consensus to be reached. In the theory of communicative action, Habermas~\citep{habermas1984theory} sets out an account on which participants in dialogue, oriented to mutual understanding, advance and redeem validity claims and move, under the unforced force of the better argument, toward a rationally motivated consensus. Shared value, on this rendering, is an agreement produced in communication oriented to agreement, and the ideal at its horizon is a consensus in which the participants have come to hold the same.

Each direction reaches something and leaves something. The similarity account is precise about the correlation between antecedent overlap and satisfaction, and it treats shared value as found and not as generated, so that it has little to say to two people whose worlds differ and who must nonetheless build something common. The consensus account is precise about the procedure by which agreement may be rationally reached, and its regulative ideal is a convergence in which difference is overcome, the participants arriving at the same view. This ideal of convergence is exactly the point at which the present account parts from it.

The relational character supplies a different conception of what shared value is. On the present account shared value is not, at its root, a pre-existing similarity that the partners discover in one another, nor a convergence in which their difference is overcome. It is a matter generated in the relation between two symbolic worlds, observable in that relation, and held in common without the two worlds becoming one. Two people may render what they most value in vocabularies that do not coincide, one in the language of duty and one in the language of care, one through a religious inheritance and one through a secular one, and the shared value they come to recognize is not the reduction of these to a single rendering. It is a matter observed in the relation between their renderings, of the kind the epistemology describes, which each continues to render in a language of their own while both come to hold it together. Shared value on this account is compatible with, and indeed rests upon, a difference the partners do not dissolve.

The generative character sets this recognition in time and marks its productivity. On the consensus picture the work of dialogue is to converge, and its completion is agreement. On the present account the relation between two differing renderings is productive of a shared value that neither party held in advance, so that the crossing generates something new held in common, beyond the mere revealing of an overlap or the negotiating of a compromise. Two people who bring genuinely different worlds into relation may generate, over years, a shared value that was in neither of their worlds at the outset and that could arise only in the relation between them. This is the deepest sense in which relational understanding does the work of building a common life, that it generates value in the relation, where the received accounts discover it in advance or split the difference between antecedent positions.

The parting from Habermas can be stated exactly, since it marks the contribution. The Habermasian ideal is a consensus in which the participants come to hold the same, and difference is a condition to be worked through toward agreement. The present account holds that a shared value can be generated and held in common precisely across a difference that is not overcome, that two worlds can hold something together while remaining two, and that the pressure toward convergence, toward the fusion of the two into one view, is the very thing the relational account warns against, since it is the point at which the recognition of shared value passes over into the effacement of one world by another. What the account keeps, against the ideal of consensus, is the difference of the two who share.
